\begin{mainquestion}{cascadedaudio}{Cascaded Audio Amplifier}

\questionpoints
Consider the audio amplifier circuit below, consisting of a common-source (CS) stage cascaded with a common-drain (CD) buffer stage. We want this amplifier to have a bandwidth spanning from 20 Hz to 20 kHz, as this roughly corresponds to the standard range of human hearing. The input signal $v_i$ has a source resistance $R_{sig} = 10\text{ k}\Omega$. 

Assume both MOSFETs are biased correctly with a transconductance of $g_m = 5\text{ mS}$. The first MOSFET has a lumped gate capacitance of $C_g = 800\text{ pF}$ (this number isn't realistic, but hey, it's just a homework problem). Assume all other internal device capacitances are negligible. Channel length modulation can be neglected ($r_o \to \infty$). 

\vspace{1em}
\begin{center}
\begin{circuitikz}[american, arrowmos]
    % Input, Source Resistance, and Bias Network
    \draw (-4,0) node[left] {$v_i$} to[R, l=$R_{sig}$, a=$10\text{ k}\Omega$, o-] (-2,0) to[C, l=$C_{in}$, -*] (0,0);
    \draw (0,0) to[R, l=$R_1$, a=$1.5\text{ M}\Omega$] (0,2.5) node[vcc] {$V_{DD}$};
    \draw (0,0) to[R, l=$R_2$, a=$1.0\text{ M}\Omega$] (0,-2.5) coordinate (gnd1) node[ground] {};
    
    % Common Source Stage
    \draw (0,0) to[short] (1.5,0) node[nmos, anchor=G] (nmos1) {};
    \draw (nmos1.S) node[ground] {};
    \draw (nmos1.D) to[R, l=$R_D$, a=$10\text{ k}\Omega$] ++(0,2) coordinate (vdd1) node[vcc] {$V_{DD}$};
    
    % Coupling Capacitor to Common Drain Stage (Straight across)
    \draw (nmos1.D) to ++(2.5,0) node[nmos, anchor=G] (nmos2) {};
    
    % Common Drain Stage
    % Align VDD of stage 2 with VDD of stage 1
    \draw (nmos2.D) -- (nmos2.D |- vdd1) node[vcc] {$V_{DD}$};
    % Align ground of RS with the input bias ground
    \draw (nmos2.S) to[R, l=$R_S$, a=$5\text{ k}\Omega$] (nmos2.S |- gnd1) node[ground] {};
    
    % Output and Load
    \draw (nmos2.S) to[short, *-o] ++(2,0) coordinate (vo) node[right] {$v_o$};
    \draw (vo) to[R, l=$R_L$, a=$1\text{ k}\Omega$] (vo |- gnd1) node[ground] {};
\end{circuitikz}
\end{center}
\vspace{1em}

% =========================================================
\begin{partquestion}{smallsignal}
{Draw the mid-band small signal equivalent circuit for \textbf{each stage individually}. For both the CS and CD stages, calculate the input resistance ($R_{in}$), output resistance ($R_o$), and open-circuit voltage gain ($A_{vo}$). \textit{Hint: At mid-band frequencies, treat large coupling capacitors as AC short circuits, and small device capacitances as AC open circuits.}}

\begin{answer}[5in]
        \textbf{Stage 1: Common Source (CS)}
        \begin{center}
        \begin{circuitikz}[scale=0.75, transform shape]
            % Fixed drawing direction to put + on top
            \draw (0,2) to[short, o-] (1,2) node[above] {$v_{gs1}$};
            \draw (0,2) to[open, v=$v_{in1}$] (0,0); 
            \draw (1,2) to[R, l=$R_{in1}$] (1,0) to[short, -o] (0,0) node[ground]{};
            \draw (3,2) to[cI, l=$g_m v_{gs1}$] (3,0);
            \draw (3,2) to[short, -*] (5,2) to[R, l=$R_D$] (5,0);
            \draw (3,0) to[short] (5,0);
            \draw (1,0) node[ground] {}; \draw (5,0) node[ground] {};
            \draw (5,2) to[short, -o] (6,2) node[right] {$v_{o1}$};
        \end{circuitikz}
        \end{center}
        
        \textbf{Input Resistance ($R_{in1}$):} Looking into the gate, the MOSFET has an infinite input impedance at mid-band frequencies. The only AC paths to ground are through the bias resistors $R_1$ (connected to AC ground at $V_{DD}$) and $R_2$ (connected to actual ground). Thus, they appear in parallel:
        \[R_{in1} = R_1 \parallel R_2 = 1.5\text{ M}\Omega \parallel 1.0\text{ M}\Omega = 600\text{ k}\Omega\]

        \textbf{Output Resistance ($R_{o1}$):} To find $R_o$, we ground the input ($v_{in1} = 0$), which makes $v_{gs1} = 0$. Consequently, the dependent current source $g_m v_{gs1}$ becomes 0 (acts as an open circuit). Looking into the drain, the only remaining path to AC ground is through $R_D$:
        \[R_{o1} = R_D = 10\text{ k}\Omega\]

        \textbf{Open-Circuit Voltage Gain ($A_{vo1}$):} The current from the dependent source $g_m v_{gs1}$ flows entirely up through $R_D$ to ground. The output voltage is $v_{o1} = -(g_m v_{gs1}) R_D$. Since $v_{in1} = v_{gs1}$, the gain is:
        \[A_{vo1} = \frac{v_{o1}}{v_{in1}} = -g_m R_D = -(5\text{ mS})(10\text{ k}\Omega) = -50\]
        
        \vspace{1em}
        \hrule
        \vspace{1em}
        
        \textbf{Stage 2: Common Drain (CD)}
        \begin{center}
        \begin{circuitikz}[scale=0.75, transform shape]
            % Corrected CD diagram to fix the grounded source error
            \draw (0,2) to[short, o-o] (2,2) node[above] {$G$};
            \draw (0,2) to[open, v=$v_{in2}$] (0,0);
            \draw (0,0) to[short, o-o] (2,0);
            \draw (1,0) node[ground] {};
            
            \draw (2,2) to[open, v=$v_{gs2}$] (4,0); 
            
            \draw (4,2) to[cI, l=$g_m v_{gs2}$] (4,0);
            \draw (4,2) node[ground, rotate=180] {}; 
            \draw (4,0) to[short, -*] (6,0) to[R, l=$R_S$] (6,-2) node[ground] {};
            \draw (6,0) to[short, -o] (7,0) node[right] {$v_{o2}$};
        \end{circuitikz}
        \end{center}

        \textbf{Input Resistance ($R_{in2}$):} Looking into the gate of the second MOSFET, there are no biasing resistors present in this simplified AC model. Because the gate oxide acts as an insulator, the resistance is infinite:
        \[R_{in2} \to \infty\]

        \textbf{Output Resistance ($R_{o2}$):} We zero the input ($v_{in2} = 0$) and apply a test voltage $v_x$ at the source. Since the gate is at AC ground, $v_{gs2} = 0 - v_x = -v_x$. The dependent source $g_m v_{gs2} = -g_m v_x$ pulls current \textit{away} from the source node, meaning looking into the source acts as a resistance of $1/g_m$ to ground. This is in parallel with the source resistor $R_S$:
        \[R_{o2} = R_S \parallel \frac{1}{g_m} = 5000\,\Omega \parallel \left(\frac{1}{0.005}\right)\,\Omega = 5000\,\Omega \parallel 200\,\Omega = \frac{1000000}{5200}\,\Omega \approx 192.3\,\Omega\]

        \textbf{Open-Circuit Voltage Gain ($A_{vo2}$):} The output voltage is created by the dependent current source flowing into $R_S$, so $v_{o2} = (g_m v_{gs2}) R_S$. The gate-to-source voltage is $v_{gs2} = v_{in2} - v_{o2}$. Substituting this into the output voltage equation gives $v_{o2} = g_m (v_{in2} - v_{o2}) R_S$. Rearranging for the ratio of $v_{o2}/v_{in2}$ gives the standard Common-Drain gain formula:
        \[A_{vo2} = \frac{v_{o2}}{v_{in2}} = \frac{g_m R_S}{1 + g_m R_S} = \frac{(5)(5)}{1 + (5)(5)} = \frac{25}{26} \approx 0.9615\]
\end{answer}

\qnewpage

% \item \textbf{(4 points)} Draw the generalized two-port equivalent model of the cascaded amplifier system, including the source and load networks. Use this model to calculate the overall mid-band voltage gain $\frac{v_o}{v_i}$.
%     \begin{answer}[3in]
%         The cascaded two-port network model is:
%         \begin{center}
%         \begin{circuitikz}[scale=0.8, transform shape]
%             % Source and Input Divider
%             \draw (0,0) to[sV, l=$v_i$] (0,2) to[R, l=$R_{sig}$] (2.5,2);
%             \draw (2.5,2) to[short, -*] (3.5,2) node[above] {$v_{in1}$} to[R, l_=$R_{in1}$] (3.5,0);
            
%             % Stage 1 Output to Stage 2 Input
%             \draw (6.5,0) to[cV, invert, l=$A_{vo1} v_{in1}$] (6.5,2) to[R, l=$R_{o1}$] (9,2);
%             \draw (9,2) to[short, -*] (10,2) node[above] {$v_{in2}$} to[open] (10,0);
%             \node at (10,1) [circle, draw, inner sep=1pt] {$\infty$}; % Indicate infinite Ri2
            
%             % Stage 2 Output to Load
%             \draw (13,0) to[cV, invert, l=$A_{vo2} v_{in2}$] (13,2) to[R, l=$R_{o2}$] (15.5,2);
%             \draw (15.5,2) to[short, -*] (16.5,2) node[right] {$v_o$} to[R, l_=$R_L$] (16.5,0);
            
%             % Grounds
%             \draw (0,0) node[ground] {}; \draw (3.5,0) node[ground] {};
%             \draw (6.5,0) node[ground] {}; \draw (10,0) node[ground] {};
%             \draw (13,0) node[ground] {}; \draw (16.5,0) node[ground] {};
            
%             % Bottom connecting lines for visual grounding
%             \draw (0,0) to[short] (3.5,0);
%             \draw (6.5,0) to[short] (10,0);
%             \draw (13,0) to[short] (16.5,0);
%         \end{circuitikz}
%         \end{center}
        
%         Using voltage dividers at the input, inter-stage, and output:
%         \[\frac{v_{in1}}{v_i} = \frac{R_{in1}}{R_{sig} + R_{in1}} = \frac{600\text{ k}\Omega}{10\text{ k}\Omega + 600\text{ k}\Omega} \approx 0.9836\]
%         \[\frac{v_{in2}}{v_{in1}} = A_{vo1} \left( \frac{R_{in2}}{R_{o1} + R_{in2}} \right) = -50 (1) = -50 \quad (\text{since } R_{in2} \to \infty)\]
%         \[\frac{v_o}{v_{in2}} = A_{vo2} \left( \frac{R_L}{R_{o2} + R_L} \right) = (0.9615) \left( \frac{1000}{192.3 + 1000} \right) = (0.9615)(0.8387) \approx 0.806\]
        
%         The overall gain is the product of these ratios:
%         \[A_v = \frac{v_o}{v_i} = (0.9836)(-50)(0.806) \approx -39.6\]
%     \end{answer}

\end{partquestion}

% =========================================================
\begin{partquestion}{lowercutoff}
{Calculate the value of $C_{in}$ to set the lower cutoff frequency $f_L = 20\text{ Hz}$. \textit{Hint: Calculate the resistance seen by the capacitor to determine the pole.}}

\begin{answer}[3.5in]
        The lower cutoff frequency is determined by the high-pass filter formed by $C_{in}$ and the resistance it sees. 
        
        Looking into the circuit from $C_{in}$, the capacitor sees $R_{sig}$ in series with the amplifier's input resistance $R_{in1}$.
        \[R_{total} = R_{sig} + R_{in1} = 10\text{ k}\Omega + 600\text{ k}\Omega = 610\text{ k}\Omega\]
        
        The high-pass pole frequency is given by:
        \[f_L = \frac{1}{2\pi R_{total} C_{in}}\]
        
        Solving for $C_{in}$:
        \[C_{in} = \frac{1}{2\pi f_L R_{total}} = \frac{1}{2\pi (20)(610 \times 10^3)}\]
        \[C_{in} = \frac{1}{24.4 \pi \times 10^6} \approx 13.05\text{ nF}\]
\end{answer}

\end{partquestion}

% =========================================================
\begin{partquestion}{uppercutoff}
{Calculate the upper cutoff frequency $f_H$ caused by the lumped gate capacitance $C_g$ interacting with the source and bias networks. Does it sufficiently cover the standard $20\text{ kHz}$ upper limit of the audio band?}

\begin{answer}[3.5in]
        The high-frequency response is dominated by the low-pass filter formed by the gate capacitance $C_g$ and the Thévenin equivalent resistance seen from the gate to ground.
        
        This equivalent resistance is the source resistance $R_{sig}$ in parallel with the input resistance $R_{in1}$:
        \[R_{th} = R_{sig} \parallel R_{in1} = 10\text{ k}\Omega \parallel 600\text{ k}\Omega\]
        \[R_{th} = \frac{(10)(600)}{10 + 600}\text{ k}\Omega = \frac{6000}{610}\text{ k}\Omega \approx 9.836\text{ k}\Omega\]
        
        The upper cutoff frequency (dominant pole) is determined by this RC time constant:
        \[f_H = \frac{1}{2\pi R_{th} C_g}\]
        \[f_H = \frac{1}{2\pi (9836\,\Omega) (800 \times 10^{-12}\text{ F})}\]
        \[f_H \approx 20,225\text{ Hz} \approx 20.2\text{ kHz}\]
        
        Yes, the upper cutoff frequency is approximately $20.2\text{ kHz}$, which is just above the $20\text{ kHz}$ requirement, meaning it perfectly covers the standard upper limit of the human audio band without significantly attenuating audible high frequencies.
\end{answer}

\end{partquestion}

\end{mainquestion}
